\documentclass[11pt]{article}
\usepackage[margin=1in]{geometry}
\usepackage{amsmath,amssymb,amsthm}
\usepackage[T1]{fontenc}
\usepackage{lmodern}
\usepackage[colorlinks=true,allcolors=blue]{hyperref}
\hypersetup{pdftitle={Loss of total-space positivity for the normalized relative Kahler-Ricci flow},pdfauthor={Alec Kriebel},pdfsubject={Unrefereed AI-assisted mathematical research note},pdfkeywords={relative Kahler-Ricci flow; positivity; projective bundles; counterexample}}
\newtheorem{theorem}{Theorem}
\newcommand{\PP}{\mathbf P}
\newcommand{\OO}{\mathcal O}
\newcommand{\MA}{\operatorname{MA}}
\newcommand{\ddbar}{\partial\bar\partial}
\title{Loss of total-space positivity for the normalized relative K\"ahler--Ricci flow}
\author{Alec Kriebel\\\small Independent researcher\\
\small ORCID: \href{https://orcid.org/0009-0001-9320-500X}{0009-0001-9320-500X}}
\date{12 October 2026 \quad (version 1.0)}
\begin{document}
\maketitle
\begin{abstract}
We give a negative answer to the total-space positivity-preservation question
for the relative K\"ahler--Ricci flow defined by Naumann on $\OO_E(r)$.
For the ample rank-two bundle $E=\OO_{\PP^1}(1)\oplus\OO_{\PP^1}(1)$,
we construct a smooth strictly positive initial metric on
$\OO_E(2)\simeq\OO_{\PP^1\times\PP^1}(2,2)$.
For the defining flow with probability-normalized fiber measures, its horizontal
curvature at one point is exactly
$1/96-(1-e^{-2t})/24$, hence equals $-1/96$ at $t=(\log2)/2$.
The argument uses the exact evolution of a second base derivative along a
stationary central fiber and an explicit global positivity certificate.
\end{abstract}

\section{Question and statement}
Naumann's report \cite{O} ends on p.~2443 with the question of preserving
positivity on the total space when a positive metric on $\OO_E(1)$ is evolved
through its $r$th power. The question remains explicit in \cite[\S3.3]{N}.
We answer it for the flow actually defined in \cite[(13)]{N}.
The flow's construction itself was already given in those sources.
The older report abbreviates the denominator without a time subscript;
we use the evolving measure explicit in the published defining equation.

Let $f:X=\PP(E^*)\to S$ use the convention $f_*\OO_E(1)=E$.
For a metric $e^{-\phi}$ on $L=\OO_E(r)$, the isomorphism
$K_{X/S}^{-1}=L\otimes f^*(\det E)^{-1}$ identifies its fiber restrictions
with anticanonical densities after a determinant frame is chosen. Define
\begin{equation}\label{flow}
 \partial_t\phi_t=\log\frac{\MA(\phi_t)}{\mu_{\phi_t}},\qquad
 \mu_\phi=\frac{e^{-\phi}\,dA_z}{I(s)},\quad
 I(s)=\int_{\PP^1}e^{-\phi(s,z)}\,dA_z.
\end{equation}
Here the fiber dimension in the example is one, and $\MA$ is curvature volume
divided by its fixed total mass. The density notation uses compatible
anticanonical frames; it is not the integral of a globally defined scalar
function against an arbitrarily chosen background area form. A determinant
frame change cancels from $\mu_\phi$. Addition of a base function changes
neither probability measure. In local coordinates, with $g=\phi_{z\bar z}>0$,
\begin{equation}\label{scalar}
 \partial_t\phi=\log g+\phi+\log I+C,
\end{equation}
where $C$ is a fixed constant determined by the area convention and the
curvature mass. In particular the reaction coefficient for this weight is one.
Omitting the fixed mass normalization in the report changes only a function
of time in the solution and leaves its curvature unchanged.

\begin{theorem}\label{thm}
For $S=\PP^1$ and $E=\OO_S(1)\oplus\OO_S(1)$, there is a smooth Hermitian
metric $e^{-\phi_0}$ on $L=\OO_E(2)$ with strictly positive curvature on the
compact total space $X$, such that its solution of \eqref{flow} satisfies
\begin{equation}\label{negative}
 (\phi_t)_{s\bar s}(0,0)=\frac1{96}-\frac{1-e^{-2t}}{24}.
\end{equation}
At $t_*=(\log2)/2$ this coefficient is $-1/96$. Fiberwise curvature remains
positive. The positive square-root metric on $\OO_E(1)$ gives admissible
initial data for the source's question.
\end{theorem}

We use the finite-time existence and smoothness statement of
\cite[Theorem~5]{N}, including smoothness on $X\times[0,\infty)$.
Its proof reduces the fiberwise flow to the normalized anticanonical
K\"ahler--Ricci flow; see also \cite{B}. No regularity or convergence of an
infinite-time limit in base directions is needed here.
Uniqueness can be checked directly: for the difference $v$ of two solutions,
$|\log I(\phi)-\log I(\widetilde\phi)|\leq\|v\|_{C^0}$ on each fiber.
At a fiber maximum of $v$ the logarithmic Hessian difference is nonpositive.
Applying the maximum principle to $v$ and $-v$ gives
an upper Dini derivative bound $D^+\|v\|_{C^0}\leq2\|v\|_{C^0}$.
Gr\"onwall's inequality proves uniqueness for common initial data and
therefore preservation of lifted symmetries.

\section{A globally positive initial metric}
Write $E=\OO_S(1)\otimes\mathbf C^2$. Then
$X\simeq\PP^1_s\times\PP^1_z$ and
$\OO_E(1)=\OO(1,1)$, so $E$ is ample and $L=\OO(2,2)$.
Set
\[
 \epsilon=\frac1{96},\qquad\delta=\frac14,\qquad
 u=|s|^2,\quad v=|z|^2,\quad p=\frac{u}{1+u},\quad x=\frac{v}{1+v}.
\]
In the finite chart take
\begin{equation}\label{initial}
 \phi_0=2\log(1+u)+\frac{2-\epsilon}{1+u}
       +2\log(1+v)+\delta\frac{u}{1+u}\left(\frac{v}{1+v}\right)^2.
\end{equation}
The logarithms are standard weights on $\OO(2,2)$ and the remaining terms
are global smooth functions. Explicitly, under $w=1/s$ subtract the
transition weight $2\log|s|^2$. The base terms become
$2\log(1+|w|^2)+(2-\epsilon)|w|^2/(1+|w|^2)$, with
$p=(1+|w|^2)^{-1}$. Under $\zeta=1/z$ subtract $2\log|z|^2$;
the fiber logarithm becomes $2\log(1+|\zeta|^2)$ and
$x^2=(1+|\zeta|^2)^{-2}$. These formulas also show smoothness in the joint
infinity chart.

In product Fubini--Study unit frames the Hessian diagonal coefficients and
squared mixed modulus are
\begin{align}
 A&=(1+u)^2(\phi_0)_{s\bar s}
     =4p+(\epsilon+\delta x^2)(1-2p),\nonumber\\
 B&=(1+v)^2(\phi_0)_{z\bar z}
     =2+2\delta p x(2-3x),\label{ABC}\\
 Q&=(1+u)^2(1+v)^2|(\phi_0)_{s\bar z}|^2
     =4\delta^2p(1-p)x^3(1-x).\nonumber
\end{align}
Direct differentiation proves these identities. They extend continuously
to the closed square $0\leq p,x\leq1$, covering every infinity point.
Throughout this square,
\[
 A\geq\epsilon+\frac{167}{48}p=:A_0,\qquad
 B\geq\frac32=:B_0,\qquad Q\leq\frac p4=:Q_0.
\]
For explicit certificates,
\begin{align*}
 A-A_0&=\delta x^2+2\delta p(1-x^2)\geq0,\\
 B-B_0&=2\delta\{(1-p)+p(1-x)(1+3x)\}\geq0,\\
 Q_0-Q&=4\delta^2p\{1-(1-p)x^3(1-x)\}\geq0.
\end{align*}
Because $A_0,B_0$ are positive,
\begin{equation}\label{det}
 AB-Q\geq A_0B_0-Q_0=\frac1{64}+\frac{159}{32}p>0.
\end{equation}
Thus the curvature is strictly positive in every direction everywhere.
This continuum certificate does not rely on sampled points.
Dividing the weight by two gives the required positive root metric.

\section{An exact second-variation equation}
Rotational symmetry in the base, preserved by uniqueness, gives on $s=0$
\begin{equation}\label{sym}
 (\phi_t)_s=(\phi_t)_{\bar s}=(\phi_t)_{s\bar z}=0.
\end{equation}
The initial central-fiber weight is
$2\log(1+|z|^2)+2-\epsilon$. Its two probability measures in
\eqref{flow} coincide, so uniqueness gives
\[
 \phi_t(0,z)=2\log(1+|z|^2)+2-\epsilon
 \quad\hbox{for every finite }t\geq0.
\]
Define
\[
 w(z,t)=(\phi_t)_{s\bar s}(0,z),\qquad
 g_F=\frac2{(1+|z|^2)^2},\qquad
 \Delta_F=\frac{(1+|z|^2)^2}{2}\partial_z\partial_{\bar z}.
\]
The sign convention makes $\Delta_F$ nonpositive on nonconstant
eigenfunctions. Differentiating \eqref{scalar}, using \eqref{sym}, yields
$(\log g)_{s\bar s}=g_F^{-1}w_{z\bar z}=\Delta_Fw$.
Smoothness and compactness justify differentiation under the fiber integral:
\[
 I_s=0,\qquad I_{s\bar s}=-I\int_{\PP^1}w\,\mu_F,
 \qquad(\log I)_{s\bar s}=-\int_{\PP^1}w\,\mu_F.
\]
Consequently the second derivative of the actual nonlinear solution obeys
the closed equation
\begin{equation}\label{jet}
 w_t=\Delta_Fw+w-\int_{\PP^1}w\,\mu_F,\qquad
 w(z,0)=\epsilon+\delta x^2.
\end{equation}
This is an exact identity, not a linearization approximation.
The geodesic curvature is
$c(\phi_t)=(\phi_t)_{s\bar s}-|(\phi_t)_{s\bar z}|^2/(\phi_t)_{z\bar z}$;
by \eqref{sym} it equals $w$ on the central fiber.

\section{Solution and strict negativity}
The pushforward of $\mu_F$ under $x=v/(1+v)$ is Lebesgue measure on
$[0,1]$, since radial integration gives $dv/(1+v)^2=dx$.
For a smooth $F(x)$,
\[
 \Delta_FF=\tfrac12\{x(1-x)F''+(1-2x)F'\}.
\]
Let $P_1=2x-1$ and $P_2=6x^2-6x+1$. Both have zero mean;
$\Delta_FP_1=-P_1$, $\Delta_FP_2=-3P_2$, and
$x^2=1/3+P_1/2+P_2/6$. Hence \eqref{jet} has solution
\begin{equation}\label{solution}
 w(x,t)=\epsilon+\frac\delta3+\frac\delta2P_1(x)
                   +\frac\delta6e^{-2t}P_2(x).
\end{equation}
To justify its identification with the actual second derivative, the difference
of two solutions of \eqref{jet} with common initial data has zero mean:
integrating its equation shows that the mean is constant. The difference
then solves $a_t=\Delta_Fa+a$; the maximum principle applied to $e^{-t}a$
proves uniqueness.

At $z=0$, $P_1=-1$ and $P_2=1$, so
$w(0,t)=\epsilon-\delta(1-e^{-2t})/6$, which is \eqref{negative}.
For $t=t_*$ it equals $-1/96$, and is negative for every
$t>\tfrac12\log(4/3)$. Because the mixed Hessian vanishes there,
this is negative curvature on the actual tangent vector $\partial/\partial s$.
All initial data were strictly positive on a compact total space, and
$t_*$ is finite. This proves Theorem~\ref{thm}.

\section{Normalization, scope, and verification}
There is a convention discrepancy between the defining equation
\cite[(13)]{N} and a literal same-weight/time reading of the displayed
$r c$ term in \cite[(14)]{N}. A calibration independent of our example is
the positive product weight
$\phi_{\rm prod}=2\log(1+|s|^2)+2\log(1+|z|^2)$.
It is stationary under \eqref{flow}. At $s=0$ its geodesic curvature is
$c=2$, the fiber Laplacian and Kodaira--Spencer terms vanish, and
$(\log I)_{s\bar s}=-2$. A same-convention $r c$ term with $r=2$
would therefore give $c_t=2$, contradicting stationarity.
Our proof uses \eqref{flow} and derives \eqref{jet} directly.
We do not assign an editorial explanation or undocumented intention to the
source. For comparison, setting $\psi=\phi/r$ and $\tau=t/r$ transforms
\eqref{scalar} into
\[
 \psi_\tau=\log\psi_{z\bar z}+r\psi+
           \log\int e^{-r\psi}\,dA_z+C',
\]
so a coefficient $r$ can occur with the corresponding weight and time
changes. Such positive rescaling leaves the occurrence of positivity loss
unchanged. Probability normalization is essential: omitting $\log I$
changes the flow and the second-variation equation.

The result resolves the source's universal positivity-preservation question
for its defining flow. It does not settle the Griffiths conjecture, exclude
special initial metrics, or address other normalizations. The example's
bundle itself already has Griffiths-positive metrics. The prior construction
of the flow is not claimed as a new contribution.

The accompanying verification package contains the manuscript source,
exact symbolic checks, analytic derivations, source and priority audit, and
adversarial review records. Symbolic computations check the algebra and
global inequality certificates; smooth parabolic existence is the cited
analytic input, not a computational conclusion. Historical first priority
is not asserted; the accompanying audit records the scope and limits of
the literature search.
In particular, Berman's preservation theorem concerns a nonnormalized
weight flow, and his earlier \cite[Remark~4.11]{B} already identifies
normalized anticanonical negativity for effectively parametrized Fano
families. Those statements do not supply this locally isotrivial
projective-space bundle example on $\OO_E(r)$. They preclude a broad
claim of first positivity loss for normalized relative flows.
A relevant Naumann habilitation listed by his university in 2025 was
identified, but its text was not recovered; this is a specific unresolved
priority lead recorded in the accompanying audit.

\paragraph{AI use and review status.}
AI tools were used extensively in developing, writing, reproducing, and
adversarially checking this result. Internal adversarial reviews were
performed by AI agents. This is an unrefereed preprint and has not received
human peer review; these internal checks do not constitute such review.

\begin{thebibliography}{9}
\bibitem{O} P.~Naumann, An approach to the Griffiths conjecture,
in \emph{Komplexe Analysis}, Oberwolfach Reports \textbf{14} (2017), no.~3,
2440--2443. \href{https://doi.org/10.4171/OWR/2017/39}{doi:10.4171/OWR/2017/39}.
\bibitem{N} P.~Naumann, An approach to the Griffiths conjecture,
\emph{Mathematical Research Letters} \textbf{28} (2021), no.~5, 1505--1523.
\href{https://doi.org/10.4310/MRL.2021.v28.n5.a10}{doi:10.4310/MRL.2021.v28.n5.a10}.
Earlier version: \href{https://arxiv.org/abs/1710.10034}{arXiv:1710.10034}.
\bibitem{B} R.~J.~Berman, Relative K\"ahler--Ricci flows and their quantization,
\emph{Analysis \& PDE} \textbf{6} (2013), no.~1, 131--180.
\href{https://doi.org/10.2140/apde.2013.6.131}{doi:10.2140/apde.2013.6.131}.
\end{thebibliography}
\end{document}
